\scanpage{76}\mspage{75}
Inversement, supposons que $F'\ll F$ et $G'\ll G$ soient tels que $F'_K\mathrel{>\mkern-14mu<}G'_K$ \deleted{\unsure{soient|soit}, désignons par $H'_{\unsure{\mathord{.}|\mathord{'}}}$} \inserted{posons}
\[
\begin{aligned}
H'_K&=F'_K\vee G'_K &&\text{on aura }H'_K\ll K\\[-.3ex]
&&&\hspace{3em}\text{par At 4}\\
K'&=\underbrace{F'_K\cap G'_K}_{\unsure{\mathord{?}|\mathord{\&}}}\deleted{=F'\cap G'}
\end{aligned}
\]

Considérons $F'$ et $H'_K$ \deleted{\unsure{on voit} aussitôt que} \inserted{\deleted{deux \unsure{raffinements|raffinement} de}} \deleted{\unsure{$F$}}raffinement de $F$ l'un, de $K$ l'autre, et tels que $F'_K\overset{\unsure{\text{défn}|\text{dfn}}}{=}F'|K$ soit compatible avec $H'_K$. Si on peut en conclure qu'ils sont compatibles, on pourrait construire leur \unsure{Sup} $H'_F$ et $H'_F\longrightarrow F$, et symétriquement $H'_G$ (\unsure{Sup} de $G'$ et de $H'_K$). On vérifierait sans mal $H'_F|K=H'_K$, et de $=$ $H'_G|K=H'_K$, donc par l'axiome de recollement, on \unsure{saurait} $H'_F\mathrel{>\mkern-14mu<}H'_G$ et on pourrait construire $H'=H'_F\vee H'_G$. \deleted{\unsure{D}}En\unsure{.} résumé, pour vérifier la partie non triviale de la proposition, on est ramené au cas particulier où
\diagram{\begin{tikzcd}[ampersand replacement=\&,row sep=1.25em,column sep=0.25em]
G \arrow[d,phantom,"{\rotatebox{90}{$\ll$}}" description] \& \deleted{\unsure{\leq}}\unsure{\subset|<} \& F \arrow[d,phantom,"{\rotatebox{90}{$\ll$}}" description] \& \text{et où}\quad F'_G\leq G' \\
G' \& \& F' \&
\end{tikzcd}}
